
% ============================================================================
\clearpage
\section*{v4.10.0-r26 Source-Only: A Duplication Rule for the Shift Tower above HS(6), and the Admissibility of HS(7) and HS(8)}
\addcontentsline{toc}{section}{v4.10.0-r26 Source-Only: A Duplication Rule for the Shift Tower above HS(6), and the Admissibility of HS(7) and HS(8)}
\label{sec:v4100-r26-hs-duplication}
% ============================================================================

\paragraph{Purpose.}
v4.10.0-r25 used two and four Dixon copies as HS(7) and HS(8) and left their admissibility open.  The atlas (v470/32, layer r58) defines
the carrier clock \(C_n=2^n\) for every \(n\) but reads only HS(0) to HS(6); it sets the HS(4) meta-septum gate and the HS(5) front/back
duplication with its cross-side debt, recommends an image audit of the HS(6) coding map before any new shift, and uses the Dixon carrier
as a route code for OP5 with an explicit refusal of Standard-Model claims.  On the author's decision (2026-11-16) this layer declares a
duplication rule for the stages above HS(6) and checks HS(7) and HS(8) against the gate and the debt register.  Exploration, not blind.
The notes are \texttt{notes/\allowbreak research/\allowbreak v4\_10\_0\_r26/} and
\texttt{notes/\allowbreak research/\allowbreak op3\_hs78\_admissibility\_2026\_11\_16/}.  Earlier layers are not rewritten.

\noindent\textbf{Source-only status.}
This v4.10.0-r26 layer is source-only.  It does not build a PDF, does not assign a new DOI,
and does not authorize a public release.  The v5.1.0 Version DOI
\href{https://doi.org/10.5281/zenodo.24237112}{10.5281/zenodo.24237112} and the Concept DOI
\href{https://doi.org/10.5281/zenodo.19210728}{10.5281/zenodo.19210728} are unchanged by it.

\subsection*{v4.10.0-r26 results}
\label{subsec:v4100-r26-results}

\paragraph{Definition v4.10.0-r26.1 (duplication rule above HS(6); a setting).}
For \(n\ge6\) put \(\mathrm{HS}(n+1):=\mathrm{HS}(n)\oplus\mathrm{HS}(n)\) with projector \(P^{(n+1)}:=\mathrm{diag}(P^{(n)},P^{(n)})=\One_2\otimes P^{(n)}\) and cross-copy
debt \(\mathrm{Debt}^{(n+1)}(T):=(\One-P^{(n+1)})\,T\,P^{(n+1)}+P^{(n+1)}\,T\,(\One-P^{(n+1)})\), in the pattern of the HS(5) duplication.  A stage is
admissible when every copy is HS(4) meta-septum closed and the debt register vanishes.  Hence HS(7) \(=2\times\)HS(6) and HS(8) \(=4\times\)HS(6).
The rule is set, not derived from the atlas.

\paragraph{Theoremlet v4.10.0-r26.2 (copy operations are free of leakage and debt).}
\begin{enumerate}[leftmargin=2em]
\item[(a)] Every operation acting only on the copy index, \(X\otimes\One\), commutes with \(\One\otimes P\); its debt and leakage vanish identically
(computed: at most \(6\cdot10^{-15}\) for two and four copies, with a placeholder projector of rank \(7\) on \(16\) channels).
\item[(b)] All constructions of r25 are copy operations and therefore admissible: Hadamard and phase on the copy qubits, the projections onto
triplet and singlet (debt at most \(3\cdot10^{-15}\)).  The triple product of r25.3 is a readout from three states, not an operator on the
carrier, and does not touch the registers.
\item[(c)] Debt arises only from motion inside a copy: an operator \(A\) with \([A,P]\ne0\) produces the same debt with or without copy mixing
(\(31.00\) in both cases).  The HS(4) gate must hold in each copy, as it holds on each side of HS(5).
\item[(d)] The stages stay nested: \(P^{(8)}\) restricted to one half is \(P^{(7)}\), to one quarter \(P^{(6)}\); the halving of r25.4 is compatible.
\end{enumerate}

\paragraph{Proof sketch.}
\([X\otimes\One,\One\otimes P]=0\) gives (a) and (b); (c) and (d) by the block structure.  Checked in \texttt{hs78\_\allowbreak admissibility.py}.  \(\square\)

\paragraph{Remark v4.10.0-r26.3 (what admissibility does and does not give).}
Under Definition r26.1, HS(7) and HS(8) are admissible extensions, and the copy qubit, the triplet--singlet split and the triple product of
r25 keep the leakage and debt registers at zero.  Two limits stand.  First, the rule is a setting; the atlas defines only HS(5) as a
duplication and recommends the image audit of HS(6) first, which is not done here.  Second, the atlas uses the shift tower as an OP5 route
grammar without particle claims, while r25 reads the copies as carriers of generations.  Admissibility shows that the two readings do not
contradict each other; it does not support the physical reading.  For r25 only the copies of \(\C\otimes\Hh\otimes\Oo\) matter, not the
atlas tower.

\paragraph{What this does not give.}
No derivation of the duplication rule, no image audit of HS(6), no statement about OP5 routes on HS(7) or HS(8), no support for the
generation reading of r25 beyond consistency.  OP3 and OP5 stay open.

\subsection*{v4.10.0-r26 status in the three frames}
\label{subsec:v4100-r26-frames}

\begin{itemize}[leftmargin=2em]
\item M, T, P: the duplication rule and the registers are frame-free; the front/back sides of HS(5) are not frames.
\end{itemize}

\subsection*{v4.10.0-r26 register entries}
\label{subsec:v4100-r26-register}

\begin{center}
\small
\begin{tabular}{@{}>{\raggedright\arraybackslash}p{0.34\textwidth}>{\raggedright\arraybackslash}p{0.58\textwidth}@{}}
\toprule
statement & status after v4.10.0-r26\\
\midrule
shift tower above HS(6) & duplication rule declared as a setting (r26.1)\\
HS(7), HS(8) & admissible under r26.1; copy operations free of leakage and debt (r26.2)\\
constructions of r25 & admissible under r26.1; physical reading not supported by it (r26.3)\\
image audit of HS(6) (atlas) & open, unchanged\\
OP3, OP5 & open\\
\bottomrule
\end{tabular}
\end{center}

\subsection*{v4.10.0-r26 ledger and guard update}
\label{subsec:v4100-r26-ledger-guard}

\[
R_{\rm chk}^{v4100\text{-}r26}=(1,1,0)
\]
for (rule declared; block computations checked in the cited script; derivation of the rule and image audit open).

The layer spends no guard:
\[
R_{\rm glob}^{(v4100\text{-}r26,+)}=(0,1,1,1,1),
\qquad
\Gamma_{\rm guard}^{v4100\text{-}r26}=(1,1,1,1,1),
\qquad
\beta_4=1/256\ \text{unspent}.
\]
The layer does not imply exact finite \(C5\) zero closure, exact all-readout terminality, \(C6/T2\), a global photon-clock/RH theorem,
or global all-readout closure.
